A near-zero cross-dataset segmentation score does not, by itself, separate visual domain shift from absence of source annotations assigned to the evaluated output class. We examine this ambiguity through a raw-mask coverage audit and fixed-budget source-mixture experiments under a declared seven-leaf operational crosswalk, with OpenEarthMap rangeland treated as a vegetation proxy whose native scope also includes shrubs. In the primary LoveDA--OpenEarthMap case, adding the tested FLAIR herbaceous source mixture increases rangeland IoU by 0.3118 but decreases mIoU over six LoveDA-supported leaves by 0.1011. A reanalysis of archived native ResNet-34 U-Net records on the same target shows the same direction: rangeland recovery of 0.2828 and supported-leaf mIoU change of -0.1304. A reversed source mixture recovers cropland with another supported-leaf cost, while Semantic Drone supplies a separate target replay of the primary CLIP checkpoints. These results extend established label-space and negative-transfer concerns through a specific audited empirical design; they introduce no new segmentation algorithm or recovery metric. Source mixtures change imagery and labels jointly, and historical checkpoint and schedule provenance remains incomplete. The evidence therefore supports bounded descriptive contrasts rather than a label-only causal effect or a successful mitigation method. Positive source-label coverage, missing-leaf recovery and supported-leaf retention should be reported together.
